% ================================================================
%  Section VII-A: KEM Primitive Evaluation
%  Data: v2-1.8ghz/raw_data/raw/kem/*.json — 100 iterations
%  Platform: RPi4 @ 1.8 GHz, INA219 1 kHz sampling
% ================================================================

\subsection{KEM Primitive Evaluation}\label{sec:kem}

We evaluate nine post-quantum key encapsulation mechanisms spanning
three families---lattice-based (ML-KEM-512/768/1024), code-based
Hamming quasi-cyclic (HQC-128/192/256), and code-based Goppa
(Classic McEliece-348864/460896/6960119)---across NIST security
levels~1, 3, and~5.  Each algorithm executes 100 iterations on the
RPi4 testbed (ARM Cortex-A72 @ \mbox{1.8\,GHz}) with board-level
energy captured via INA219 at \mbox{1\,kHz}.
Table~\ref{tab:kem_full} presents the complete benchmark; the
$\rho$~column normalises total cycle time to ML-KEM at the same
NIST level.

%% ---------------------------------------------------------------
%%  TABLE: Comprehensive KEM Benchmark
%% ---------------------------------------------------------------
\begin{table*}[t]
\centering
\caption{Comprehensive KEM primitive benchmark on RPi4 (100 iterations, INA219 1\,kHz).
  ML-KEM completes a full keygen--encaps--decaps cycle in $<$0.72\,ms
  and $<$2.5\,mJ at every NIST level; alternatives are 255--12\,823$\times$ slower.}
\label{tab:kem_full}
\footnotesize
\setlength{\tabcolsep}{3pt}
\begin{tabular}{@{}ll rrr r rrr r rrr r c@{}}
\toprule
& & \multicolumn{3}{c}{\textbf{Latency (ms)}}
& \textbf{$T_\text{tot}$}
& \multicolumn{3}{c}{\textbf{Energy ($\boldsymbol{\mu}$J)}}
& \textbf{$E_\text{tot}$}
& \multicolumn{3}{c}{\textbf{Size (B)}}
& \textbf{$\boldsymbol{\rho}$}
& \textbf{UAV} \\
\cmidrule(lr){3-5}\cmidrule(lr){7-9}\cmidrule(lr){11-13}
\textbf{Algorithm} & \textbf{L}
 & \textbf{Kg} & \textbf{Enc} & \textbf{Dec}
 & \textbf{(ms)}
 & \textbf{$E_\text{kg}$} & \textbf{$E_\text{enc}$} & \textbf{$E_\text{dec}$}
 & \textbf{($\mu$J)}
 & \textbf{PK} & \textbf{CT} & \textbf{SK}
 &
 & \\
\midrule
ML-KEM-512   & 1 & 0.240 & 0.152 & 0.158 & 0.550
  &     830 &     526 &     551 &   1\,907
  &    800 &    768 & 1\,632
  & 1.0$\times$ & $\checkmark\!\checkmark\!\checkmark$ \\
ML-KEM-768   & 3 & 0.237 & 0.182 & 0.185 & 0.604
  &     823 &     636 &     653 &   2\,112
  &  1\,184 & 1\,088 & 2\,400
  & 1.0$\times$ & $\checkmark\!\checkmark\!\checkmark$ \\
ML-KEM-1024  & 5 & 0.273 & 0.216 & 0.231 & 0.720
  &     947 &     754 &     806 &   2\,507
  &  1\,568 & 1\,568 & 3\,168
  & 1.0$\times$ & $\checkmark\!\checkmark\!\checkmark$ \\
\addlinespace
HQC-128      & 1 & 22.33 & 44.68 & 73.39  & 140.4
  & 83\,153 & 167\,590 & 284\,510 & 535\,253
  & 2\,249 & 4\,433 & ---
  & 255$\times$ & $\checkmark$ \\
HQC-192      & 3 & 67.68 & 135.8 & 211.9  & 415.4
  & 257\,928 & 537\,757 & 856\,128 & 1\,651\,813
  & 4\,522 & 8\,978 & ---
  & 688$\times$ & $\sim$ \\
HQC-256      & 5 & 124.4 & 250.2 & 392.6  & 767.2
  & 492\,250 & 1\,023\,984 & 1\,620\,633 & 3\,136\,867
  & 7\,245 & 14\,421 & ---
  & 1\,066$\times$ & $\sim$ \\
\addlinespace
McEliece-348 & 1 & 351.0 & 0.466 & 55.98  & 407.4
  & 1\,527\,273 & 1\,619 & 208\,294 & 1\,737\,186
  & 261\,120 &    96 & ---
  & 741$\times$ & $\times^{\dagger}$ \\
McEliece-460 & 3 & 1\,281 & 0.872 & 89.59 & 1\,371
  & 5\,682\,343 & 3\,135 & 339\,956 & 6\,025\,434
  & 524\,160 &   156 & ---
  & 2\,272$\times$ & $\times^{\dagger}$ \\
McEliece-8192& 5 & 9\,021 & 1.894 & 210.1 & 9\,233
  & 39\,746\,860 & 6\,849 & 829\,782 & 40\,583\,491
  & 1\,357\,824 &   208 & ---
  & 12\,823$\times$ & $\times$ \\
\bottomrule
\end{tabular}
\begin{tablenotes}
\item $T_\text{tot}=T_\text{kg}+T_\text{enc}+T_\text{dec}$;\;
  $\rho$: total cycle time relative to ML-KEM at the same NIST level.\;
  $\dagger$\,PK sizes (261\,KB--1.33\,MB) prohibitive for constrained UAV links.\;
  UAV: $\checkmark\!\checkmark\!\checkmark$\,=\,real-time viable,
  $\checkmark$\,=\,viable with extended interval,
  $\sim$\,=\,marginal,
  $\times$\,=\,non-viable.
\end{tablenotes}
\end{table*}

\subsubsection{Latency Analysis}

ML-KEM operates sub-millisecond across all three security levels:
keygen completes in 0.237--0.273\,ms, encapsulation in
0.152--0.216\,ms, and decapsulation in 0.158--0.231\,ms,
yielding full-cycle times of 0.550--0.720\,ms.
The three operations are near-balanced---keygen accounts for
38--44\% of the cycle, encapsulation 28--30\%, and decapsulation
29--32\%---confirming that no single phase constitutes a
bottleneck on ARM Cortex-A72 hardware.

HQC exhibits a sharply different profile.
At Level~1, full-cycle latency is 140.4\,ms ($\rho = 255\times$),
rising to 415.4\,ms at Level~3 ($688\times$) and 767.2\,ms
at Level~5 ($1{,}066\times$).
Decapsulation dominates every HQC variant, consuming
51--52\% of the cycle due to the cost of syndrome decoding
over quasi-cyclic codes; keygen is the lightest phase at only~16\%.

Classic McEliece reveals an extreme keygen asymmetry.
McEliece-348864 requires 351.0\,ms for key generation (86\% of
its 407.4\,ms cycle), while McEliece-8192128 demands 9.021\,s
(98\% of its 9.233\,s cycle) due to probabilistic Goppa-code
generation with a heavy-tail distribution
($\sigma/\mu = 74\%$ at Level~5).
Encapsulation, however, is remarkably fast---0.466--1.894\,ms,
within an order of magnitude of ML-KEM---but this single fast
operation cannot offset the multi-second keygen penalty.
Fig.~\ref{fig:kem_three_panel} visualises these per-operation
contrasts on a logarithmic scale.

%% ---------------------------------------------------------------
%%  FIGURE 1: Three-panel horizontal bar chart (Keygen|Encaps|Decaps)
%% ---------------------------------------------------------------
\begin{figure*}[t]
\centering
\begin{subfigure}[b]{0.32\textwidth}
\centering
\begin{tikzpicture}
\begin{axis}[
  xbar,
  width=\textwidth,
  height=5.2cm,
  bar width=5.5pt,
  xlabel={Keygen (ms)},
  xmode=log, log basis x=10,
  xmin=0.1, xmax=20000,
  symbolic y coords={ML512,ML768,ML1024,HQ128,HQ192,HQ256,MC348,MC460,MC8192},
  ytick=data,
  y tick label style={font=\scriptsize},
  x tick label style={font=\scriptsize},
  xlabel style={font=\scriptsize},
  nodes near coords,
  nodes near coords style={font=\tiny, anchor=west},
  point meta=explicit symbolic,
  enlarge y limits=0.08,
  grid=major, grid style={dotted, gray!40},
]
\addplot[fill=ptblue!45, draw=ptblue!85!black] coordinates {
  (0.240,ML512)  [0.24]
  (0.237,ML768)  [0.24]
  (0.273,ML1024) [0.27]
  (22.33,HQ128)  [22.3]
  (67.68,HQ192)  [67.7]
  (124.4,HQ256)  [124]
  (351.0,MC348)  [351]
  (1281,MC460)   [1281]
  (9021,MC8192)  [9021]
};
\end{axis}
\end{tikzpicture}
\caption{Keygen}\label{fig:kem_tp_kg}
\end{subfigure}
\hfill
\begin{subfigure}[b]{0.32\textwidth}
\centering
\begin{tikzpicture}
\begin{axis}[
  xbar,
  width=\textwidth,
  height=5.2cm,
  bar width=5.5pt,
  xlabel={Encaps (ms)},
  xmode=log, log basis x=10,
  xmin=0.05, xmax=1000,
  symbolic y coords={ML512,ML768,ML1024,HQ128,HQ192,HQ256,MC348,MC460,MC8192},
  ytick=data,
  y tick label style={font=\scriptsize},
  x tick label style={font=\scriptsize},
  xlabel style={font=\scriptsize},
  nodes near coords,
  nodes near coords style={font=\tiny, anchor=west},
  point meta=explicit symbolic,
  enlarge y limits=0.08,
  grid=major, grid style={dotted, gray!40},
]
\addplot[fill=ptgold!50, draw=ptgold!85!black] coordinates {
  (0.152,ML512)  [0.15]
  (0.182,ML768)  [0.18]
  (0.216,ML1024) [0.22]
  (44.68,HQ128)  [44.7]
  (135.8,HQ192)  [136]
  (250.2,HQ256)  [250]
  (0.466,MC348)  [0.47]
  (0.872,MC460)  [0.87]
  (1.894,MC8192) [1.89]
};
\end{axis}
\end{tikzpicture}
\caption{Encaps}\label{fig:kem_tp_enc}
\end{subfigure}
\hfill
\begin{subfigure}[b]{0.32\textwidth}
\centering
\begin{tikzpicture}
\begin{axis}[
  xbar,
  width=\textwidth,
  height=5.2cm,
  bar width=5.5pt,
  xlabel={Decaps (ms)},
  xmode=log, log basis x=10,
  xmin=0.05, xmax=1000,
  symbolic y coords={ML512,ML768,ML1024,HQ128,HQ192,HQ256,MC348,MC460,MC8192},
  ytick=data,
  y tick label style={font=\scriptsize},
  x tick label style={font=\scriptsize},
  xlabel style={font=\scriptsize},
  nodes near coords,
  nodes near coords style={font=\tiny, anchor=west},
  point meta=explicit symbolic,
  enlarge y limits=0.08,
  grid=major, grid style={dotted, gray!40},
]
\addplot[fill=ptrose!45, draw=ptrose!85!black] coordinates {
  (0.158,ML512)  [0.16]
  (0.185,ML768)  [0.19]
  (0.231,ML1024) [0.23]
  (73.39,HQ128)  [73.4]
  (211.9,HQ192)  [212]
  (392.6,HQ256)  [393]
  (55.98,MC348)  [56.0]
  (89.59,MC460)  [89.6]
  (210.1,MC8192) [210]
};
\end{axis}
\end{tikzpicture}
\caption{Decaps}\label{fig:kem_tp_dec}
\end{subfigure}
\caption{Three-panel KEM latency comparison (log scale).
  ML-KEM remains sub-millisecond for all operations ($<$0.28\,ms);
  McEliece encaps (0.47--1.89\,ms) rivals ML-KEM yet its keygen
  spans 0.35--9.02\,s.  HQC follows a consistent
  Kg\,$<$\,Enc\,$<$\,Dec ordering with decaps consuming 51--52\%
  of each cycle.}
\label{fig:kem_three_panel}
\end{figure*}

\subsubsection{Energy Analysis}

Energy measurements track the same structural patterns but amplify
the gaps because power draw is not constant: McEliece keygen
sustains 4.28--4.40\,W (0.96--1.08\,W above the 3.32\,W idle
baseline), whereas ML-KEM operates at 3.44--3.52\,W, barely
above idle.
The total energy for one ML-KEM cycle ranges from 1.907\,mJ
(Level~1) to 2.507\,mJ (Level~5)---comfortably within the
sub-second rekey budget of any battery-powered UAV.
HQC consumes 535\,mJ at Level~1, rising to 1.652\,J and 3.137\,J
at Levels~3 and~5 respectively; decapsulation alone accounts for
284--1\,621\,mJ, i.e.\ 51--52\% of total cycle energy.
McEliece energy is dominated almost entirely by keygen: 1.527\,J
(88\%) at Level~1, 5.682\,J (94\%) at Level~3, and 39.747\,J
(98\%) at Level~5.  The Level~5 total of 40.583\,J is
$16{,}189\times$ that of ML-KEM-1024.
Fig.~\ref{fig:kem_energy_stacked} decomposes these per-operation
contributions.

%% ---------------------------------------------------------------
%%  FIGURE 2: Stacked energy breakdown
%% ---------------------------------------------------------------
\begin{figure}[t]
\centering
\begin{tikzpicture}
\begin{axis}[
  xbar stacked,
  width=\columnwidth,
  height=5.8cm,
  xmode=log, log basis x=10,
  xlabel={Energy ($\mu$J) --- log scale},
  xlabel style={font=\scriptsize},
  x tick label style={font=\scriptsize},
  symbolic y coords={MC8192,MC460,MC348,HQ256,HQ192,HQ128,ML1024,ML768,ML512},
  ytick=data,
  y tick label style={font=\scriptsize},
  legend style={at={(0.98,0.02)}, anchor=south east, font=\scriptsize,
    legend columns=3},
  bar width=5.5pt,
  xmin=100,
  enlarge y limits=0.08,
  every axis plot/.append style={fill opacity=0.85},
  axis x line*=bottom,
  axis y line*=left,
]
\addplot[fill=ptblue!40, draw=ptblue!80!black] coordinates {
  (830,ML512) (823,ML768) (947,ML1024)
  (83153,HQ128) (257928,HQ192) (492250,HQ256)
  (1527273,MC348) (5682343,MC460) (39746860,MC8192)
};
\addplot[fill=ptgold!50, draw=ptgold!80!black] coordinates {
  (526,ML512) (636,ML768) (754,ML1024)
  (167590,HQ128) (537757,HQ192) (1023984,HQ256)
  (1619,MC348) (3135,MC460) (6849,MC8192)
};
\addplot[fill=ptrose!40, draw=ptrose!80!black] coordinates {
  (551,ML512) (653,ML768) (806,ML1024)
  (284510,HQ128) (856128,HQ192) (1620633,HQ256)
  (208294,MC348) (339956,MC460) (829782,MC8192)
};
\legend{Keygen,Encaps,Decaps}
\end{axis}
\end{tikzpicture}
\caption{Stacked energy breakdown by KEM operation (log scale).
  ML-KEM total energy is $\sim$\!2\,mJ across all levels.
  McEliece keygen dominates ($>$97\% at L5) while HQC
  decaps is the heaviest component ($>$50\% at every level).}
\label{fig:kem_energy_stacked}
\end{figure}

\subsubsection{Bandwidth--Energy Trade-off}

The UAV link budget constrains both computation and bandwidth.
Fig.~\ref{fig:kem_pareto} plots total bandwidth cost (PK\,+\,CT)
against total cycle energy for all nine variants.
ML-KEM occupies the Pareto-optimal lower-left corner: 1.568--3.136\,KB
total bandwidth and 1.907--2.507\,mJ energy.  Scaling from Level~1
to Level~5 costs only $1.96\times$ in PK size and $1.31\times$ in
energy---the most graceful scaling of any family.

HQC bandwidth is moderate (6.7--21.7\,KB) but energy is
281--1{,}251$\times$ that of ML-KEM.
McEliece offers the smallest ciphertext of any family
(96--208\,B), yet its public keys are 261\,KB at Level~1,
524\,KB at Level~3, and 1.33\,MB at Level~5---pushing it to
the far right of the plot.  At a typical UAV telemetry rate
of 250\,kbps, transmitting the McEliece-8192128 public key alone
requires 43.5\,s, consuming the entire rekey window.

%% ---------------------------------------------------------------
%%  FIGURE 3: Pareto bubble — Bandwidth vs Energy
%% ---------------------------------------------------------------
\begin{figure}[t]
\centering
\begin{tikzpicture}
\begin{axis}[
  width=\columnwidth,
  height=6.2cm,
  xlabel={PK + CT (bytes) --- log scale},
  ylabel={Total Energy ($\mu$J) --- log scale},
  xlabel style={font=\scriptsize},
  ylabel style={font=\scriptsize},
  x tick label style={font=\scriptsize},
  y tick label style={font=\scriptsize},
  xmode=log, ymode=log,
  log basis x=10, log basis y=10,
  xmin=500, xmax=2000000,
  ymin=500, ymax=100000000,
  grid=major, grid style={dotted, gray!40},
  legend style={at={(0.02,0.98)}, anchor=north west, font=\scriptsize},
]
% ML-KEM (PK+CT): 1568, 2272, 3136
\addplot[only marks, mark=*, color=ptblue, mark size=4.5pt,
  fill=ptblue!40] coordinates {
  (1568, 1907) (2272, 2112) (3136, 2507)
};
\addlegendentry{ML-KEM}
% HQC (PK+CT): 6682, 13500, 21666
\addplot[only marks, mark=triangle*, color=ptgold, mark size=4.5pt,
  fill=ptgold!50] coordinates {
  (6682, 535253) (13500, 1651813) (21666, 3136867)
};
\addlegendentry{HQC}
% McEliece (PK+CT): 261216, 524316, 1358032
\addplot[only marks, mark=square*, color=ptrose, mark size=4.5pt,
  fill=ptrose!40] coordinates {
  (261216, 1737186) (524316, 6025434) (1358032, 40583491)
};
\addlegendentry{McEliece}

% Data labels
\node[font=\tiny, anchor=south, text=ptblue!80!black]
  at (axis cs:1568,1907) {512};
\node[font=\tiny, anchor=south west, text=ptblue!80!black]
  at (axis cs:2272,2112) {768};
\node[font=\tiny, anchor=south west, text=ptblue!80!black]
  at (axis cs:3136,2507) {1024};
\node[font=\tiny, anchor=south west, text=ptgold!80!black]
  at (axis cs:6682,535253) {128};
\node[font=\tiny, anchor=south west, text=ptgold!80!black]
  at (axis cs:13500,1651813) {192};
\node[font=\tiny, anchor=south west, text=ptgold!80!black]
  at (axis cs:21666,3136867) {256};
\node[font=\tiny, anchor=south, text=ptrose!80!black]
  at (axis cs:261216,1737186) {348};
\node[font=\tiny, anchor=south, text=ptrose!80!black]
  at (axis cs:524316,6025434) {460};
\node[font=\tiny, anchor=south, text=ptrose!80!black]
  at (axis cs:1358032,40583491) {8192};

% Pareto front annotation
\draw[dashed, ptblue!60, thick]
  (axis cs:1568,1907) -- (axis cs:3136,2507);
\node[font=\tiny, text=ptblue!70!black, rotate=8, anchor=south]
  at (axis cs:2200,1650) {Pareto front};
\end{axis}
\end{tikzpicture}
\caption{Bandwidth (PK+CT) vs.\ total cycle energy (log--log).
  ML-KEM dominates the Pareto front: 1.5--3.1\,KB bandwidth and
  $<$2.5\,mJ energy.  McEliece public keys ($>$261\,KB) push it
  far right despite having the smallest CT (96--208\,B).}
\label{fig:kem_pareto}
\end{figure}

\subsubsection{NIST Level Scaling}

Fig.~\ref{fig:kem_level_scaling} isolates the scaling behaviour
from Level~1 to Level~5 in two panels: combined encapsulation +
decapsulation latency and public-key size.
ML-KEM scales most gracefully: Enc\,+\,Dec grows from 0.310\,ms to
0.447\,ms ($1.44\times$), and PK from 800\,B to 1\,568\,B
($1.96\times$)---both below $2\times$.
HQC latency scales $5.4\times$ (118.1\,ms $\to$ 642.8\,ms),
with PK growing $3.2\times$ (2\,249\,B $\to$ 7\,245\,B).
McEliece PK size explodes $5.2\times$ (261\,KB $\to$ 1.33\,MB),
while its Enc\,+\,Dec rises $3.8\times$ (56.4\,ms $\to$
212.0\,ms)---driven entirely by the decapsulation growth, since
encapsulation only moves from 0.466 to 1.894\,ms.

These scaling factors have direct implications for UAV rekey
interval selection: ML-KEM can upgrade from Level~1 to Level~5
with negligible impact ($+$0.17\,ms per cycle), whereas HQC and
McEliece force significant interval relaxation or pre-provisioning
strategies at higher levels.

%% ---------------------------------------------------------------
%%  FIGURE 4: NIST Level Scaling — Dual Panel
%% ---------------------------------------------------------------
\begin{figure}[t]
\centering
\begin{subfigure}[b]{\columnwidth}
\centering
\begin{tikzpicture}
\begin{axis}[
  width=0.92\columnwidth,
  height=3.8cm,
  xlabel={NIST Level},
  ylabel={Enc+Dec (ms)},
  xlabel style={font=\scriptsize},
  ylabel style={font=\scriptsize},
  x tick label style={font=\scriptsize},
  y tick label style={font=\scriptsize},
  ymode=log,
  xtick={1,3,5},
  xticklabels={L1,L3,L5},
  ymin=0.1, ymax=1000,
  grid=major, grid style={dotted, gray!40},
  mark size=3.5pt,
  legend style={at={(0.02,0.98)}, anchor=north west, font=\scriptsize,
    legend columns=3},
  thick,
]
\addplot[color=ptblue, mark=square*, thick] coordinates {
  (1, 0.310) (3, 0.367) (5, 0.447)
};
\addlegendentry{ML-KEM}
\addplot[color=ptgold, mark=triangle*, thick] coordinates {
  (1, 118.07) (3, 347.7) (5, 642.8)
};
\addlegendentry{HQC}
\addplot[color=ptrose, mark=diamond*, thick] coordinates {
  (1, 56.45) (3, 90.46) (5, 212.0)
};
\addlegendentry{McEliece}
\end{axis}
\end{tikzpicture}
\caption{Enc+Dec combined latency across NIST levels}\label{fig:kem_ls_lat}
\end{subfigure}

\vspace{2mm}
\begin{subfigure}[b]{\columnwidth}
\centering
\begin{tikzpicture}
\begin{axis}[
  width=0.92\columnwidth,
  height=3.8cm,
  xlabel={NIST Level},
  ylabel={PK Size (bytes)},
  xlabel style={font=\scriptsize},
  ylabel style={font=\scriptsize},
  x tick label style={font=\scriptsize},
  y tick label style={font=\scriptsize},
  ymode=log,
  xtick={1,3,5},
  xticklabels={L1,L3,L5},
  ymin=100, ymax=5000000,
  grid=major, grid style={dotted, gray!40},
  mark size=3.5pt,
  legend style={at={(0.98,0.02)}, anchor=south east, font=\scriptsize,
    legend columns=3},
  thick,
]
\addplot[color=ptblue, mark=square*, thick] coordinates {
  (1, 800) (3, 1184) (5, 1568)
};
\addlegendentry{ML-KEM}
\addplot[color=ptgold, mark=triangle*, thick] coordinates {
  (1, 2249) (3, 4522) (5, 7245)
};
\addlegendentry{HQC}
\addplot[color=ptrose, mark=diamond*, thick] coordinates {
  (1, 261120) (3, 524160) (5, 1357824)
};
\addlegendentry{McEliece}
\end{axis}
\end{tikzpicture}
\caption{Public key size across NIST levels}\label{fig:kem_ls_pk}
\end{subfigure}
\caption{NIST level scaling.  ML-KEM scales gracefully
  ($<$2$\times$ from L1$\to$L5 in both latency and PK size).
  HQC latency grows 5.4$\times$ while McEliece PK explodes
  5.2$\times$, reaching 1.33\,MB at Level~5.}
\label{fig:kem_level_scaling}
\end{figure}

\subsubsection{McEliece Encapsulation Advantage}

Despite its prohibitive keygen and PK sizes, McEliece possesses
one noteworthy property: ultrafast encapsulation.
Fig.~\ref{fig:kem_encaps_ct} overlays encapsulation latency (bars)
with ciphertext size (diamonds) on dual axes.
McEliece encaps completes in 0.466--1.894\,ms---only
$3.1\times$--$8.8\times$ slower than ML-KEM and
$96\times$--$132\times$ faster than HQC at matched levels.
McEliece also produces the smallest ciphertexts: 96\,B at Level~1,
156\,B at Level~3, and 208\,B at Level~5, compared to ML-KEM's
768--1\,568\,B and HQC's 4\,433--14\,421\,B.
However, this encapsulation advantage is entirely negated by the
keygen and PK-transmission costs.  A full KEM exchange requires
shipping the public key, and even McEliece-348864's 261\,KB PK
takes 8.4\,s to transmit at 250\,kbps---orders of magnitude longer
than its 0.466\,ms encaps.

%% ---------------------------------------------------------------
%%  FIGURE 5: Encapsulation speed + CT size dual-axis
%% ---------------------------------------------------------------
\begin{figure}[t]
\centering
\begin{tikzpicture}
\begin{axis}[
  ybar,
  width=\columnwidth,
  height=5.2cm,
  bar width=7pt,
  ylabel={Encaps Latency (ms)},
  ylabel style={font=\scriptsize},
  y tick label style={font=\scriptsize},
  ymode=log, log basis y=10,
  ymin=0.05, ymax=500,
  symbolic x coords={ML512,ML768,ML1024,HQ128,HQ192,HQ256,MC348,MC460,MC8192},
  xtick=data,
  x tick label style={rotate=35, anchor=east, font=\scriptsize},
  axis y line*=left,
  legend style={at={(0.02,0.98)}, anchor=north west, font=\scriptsize},
  every axis plot/.append style={fill opacity=0.85},
  enlarge x limits=0.08,
  nodes near coords,
  nodes near coords style={font=\tiny},
]
\addplot[fill=ptgold!50, draw=ptgold!80!black] coordinates {
  (ML512,0.152) (ML768,0.182) (ML1024,0.216)
  (HQ128,44.68) (HQ192,135.8) (HQ256,250.2)
  (MC348,0.466) (MC460,0.872) (MC8192,1.894)
};
\addlegendentry{Encaps latency}
\end{axis}
\begin{axis}[
  width=\columnwidth,
  height=5.2cm,
  ylabel={Ciphertext Size (bytes)},
  ylabel style={font=\scriptsize},
  y tick label style={font=\scriptsize},
  ymode=log, log basis y=10,
  ymin=30, ymax=50000,
  symbolic x coords={ML512,ML768,ML1024,HQ128,HQ192,HQ256,MC348,MC460,MC8192},
  xtick=\empty,
  axis y line*=right,
  axis x line=none,
  enlarge x limits=0.08,
  legend style={at={(0.98,0.98)}, anchor=north east, font=\scriptsize},
]
\addplot[only marks, mark=diamond*, color=ptrose, mark size=4pt, thick]
  coordinates {
    (ML512,768) (ML768,1088) (ML1024,1568)
    (HQ128,4433) (HQ192,8978) (HQ256,14421)
    (MC348,96) (MC460,156) (MC8192,208)
  };
\addlegendentry{CT size}
\end{axis}
\end{tikzpicture}
\caption{Encapsulation latency (bars, left axis) and ciphertext size
  (diamonds, right axis).  McEliece achieves the smallest CT
  (96--208\,B) and near sub-ms encaps, competitive with ML-KEM;
  however, its massive PK (261\,KB--1.33\,MB) and multi-second
  keygen negate this advantage for UAV rekey.}
\label{fig:kem_encaps_ct}
\end{figure}

\subsubsection{UAV Viability Assessment}

Synthesising timing, energy, and bandwidth metrics yields a clear
hierarchy.  ML-KEM is the only family that achieves real-time
viability ($\checkmark\!\checkmark\!\checkmark$) at all three NIST
levels: sub-millisecond cycles, sub-2.5\,mJ energy, and
sub-3.2\,KB total bandwidth (PK\,+\,CT).  Only ML-KEM enables
sub-second periodic rekey on ARM Cortex-A72 hardware---a 30\,s
rekey interval incurs $<$0.01\% overhead.

HQC-128 is marginally viable ($\checkmark$) with a 140.4\,ms cycle
and 535\,mJ energy, but Levels~3 and~5 push total energy above 1.6\,J
and cycle time above 415\,ms, requiring extended rekey intervals for
battery conservation.  At Level~5, the 767.2\,ms cycle and 3.137\,J
energy make HQC practical only if rekey frequency is relaxed to
$\geq$60\,s.

Classic McEliece is non-viable ($\times$) at all levels due to the
compounding effects of multi-second keygen (351\,ms--9.02\,s), extreme
energy (1.74--40.6\,J), and bandwidth-prohibitive public keys
(261\,KB--1.33\,MB).  Even at Level~1, the PK transmission time
(8.4\,s at 250\,kbps) alone exceeds the computational cost of an
entire ML-KEM-1024 rekey cycle by four orders of magnitude.

These results establish ML-KEM as the unambiguous choice for the
KEM component of a post-quantum UAV tunnel; subsequent sections
evaluate signature and AEAD primitives under the same methodology.
